\sscp{To include in grammar sketches}{15-07-02}
In future grammars
and grammar sketches of languages in the region,
this study has touched on a number of areas that should be included
to help address descriptive,
typological, and comparative issues of interest and relevance.
Some topics listed below are directly related to comparative
issues in the morphosyntax
and typology that we have touched on.
Others are derivative topics that,
if addressed, are likely to contribute to a clearer comparative
picture of the morphosyntax
and typologies of the languages across the region.
(See \citealp{Payne1997,Dixon2010} for more in-depth discussion
of these issues cross-linguistically.) The list below is not
exhaustive, but includes:

\begin{itemize}
	\item	The grammar
				and typology of transitive clauses looking at NP core arguments
				separately from pronominal arguments.
				What in the grammar or morphology signals to the hearer who (Actor)
				is doing what (activity) to whom (Undergoer)? Is ergativity marked?
				Is there a different marking for nominative
				and accusative roles? Are there other patterns driving the syntax?

	\item	The grammar
				and typology of intransitive clauses,
				looking at both active (\it{activity,
				accomplishment}) and non-active (\it{state, process,
				change-of-state, resulting state}) verbs.
				Is S\sub{A} marked the same or different than S\sub{U}? In the syntax? In
				the morphology? In the choice of pronouns? Given the prevalence
				but sporadic distribution of Split-S languages in both Austronesian
				and Papuan languages in the region, this is important.

	\item	Do verbs of motion
				and posture (and some experiencer verbs) have (optionally) different
				grammar than what might be expected? Specifically,
				are some syntactically transitive,
				with Actor and Undergoer being coreferential?

	\item	Are some experiencer verbs marked as reflexives?

	\item	What is the most unmarked order in ditransitive verbs between
				Recipient (R), and Theme (T), that is,
				the thing being exchanged or transferred? What happens in the
				grammar if/when the unmarked order is changed?

	\item	Do some verbs take subject-indexing
				and others not? If there is subject-indexing on (some) verbs,
				are there different paradigms? Give the full paradigms,
				and give all paradigms found.
				For example, are there different paradigms for verbs with V-onsets
				from those with C-onsets or CC-onsets,
				such as in \tsc{Rote-Meto} languages? Are there different morphophonemic
				processes that come into play,
				such as those found in \tsc{\ili{Ambon}-Seram} languages (and \ili{Banda},
				and some STB languages)?

	\item	Are there dual
				and trial forms of pronouns? How are base pronouns modified to
				say ‘the two of us/you/them’,
				‘the three of us/you/them’,
				‘the four of us/you/them’?

	\item	Are there applicatives on the verb that repackage the role
				structure of the verb,
				such as transitivise intransitive verbs,
				verbalise nouns, increase the valence by making the verb causative,
				decrease the valence by incorporating arguments,
				indicate dative raising, etc.?

	\item	What is the normal position in a clause for standard negation?
				For other types of negation?

	\item	Is there bipartite (split,
				embracing) standard negation? Do the younger generations of speakers
				continue to use it?

	\item	What slots in the clause can various TAM markers take? Are
				there patterned differences between pre- or post-verbal auxiliaries
				and clause-final auxiliaries? Is reduplication involved with
				indicating some aspects?
				
	\item	How are various non-verbal predicates indicated? Can a subset
				of TAM markers be used with non-verbal predicates?

	\item	Is there a copula,
				or more likely a secondary sense of a verb or adposition used
				with a copula-like function?

	\item	How is phrasal possession indicated? How is predicative possession
				indicated? Provide full paradigms for all types of possession.
				(See \srf{02-03-01-01}.)
	
	\item	Is there more than one construction for phrasal possession?
				If so, what is the difference? Is it alienable/inalienable noun classes? Bound vs.
				free roots? How are the two (or more) constructions marked? When
				there is an overt possessor,
				is it a preposed possessor,
				a postposed possessor, or can there be both? Check several body
				parts (\it{head, hand/arm, foot/leg}), close kin terms (\it{mother,
				father, child}), general things (\it{house, machete, boat, shirt}).
				Give full words, not just roots (\it{lima-n,} not \it{lima-}).

	\item	Check noun + noun part-whole relationships such as: \it{pig’s
				head, pig’s tail, roof of house, door of house,
				tree root, tree branch,
				tree trunk, tree leaf.}

	\item	Is there crossover? Can some nouns be used with more than one
				phrasal possessive construction? Is there a difference in meaning
				when that happens?
				
	\item	Is there a kind of zero-marking on nouns when they are talked
				about abstractly, in general, dismembered, or when their part-whole association
				is well established in the discourse? (E.g.
				\ili{Buru} \it{isi-t} is meat dried for taking to market; vs.,
				\it{isi-n} ‘its meat’.)
				
	\item	Are there vocative markers that are the same as,
				or similar to certain possessive markers (particularly with close
				kin terms)?
				
	\item	Is there a human/non-human,
				or animate/inanimate distinction in third person pronouns or
				possessive constructions?
				
	\item	What is the order in the NP of Head + Attribute? Head + Number/Quantifier?
				Head + Deictic? Head + Relative Clause? In natural text,
				can they, or how do they concatenate multiple attributes? How
				do they string together Head + Attribute + Number/Quantifier
				+ Relative Clause + Deictic? Are there limitations?

	\item	What is the basic system of spatial/temporal/referential deixis?
				Does it interact with directional terms? Are there other complexities?

	\item	What connectors are used to link clauses
				and sentences and paragraphs? These mechanisms are usually much
				more than just the word class of conjunctions alone.

	\item	In natural (recorded,
				not written) discourse, how are presentational clauses structured
				to introduce main participants with on-going relevance into the
				narrative?
				
	\item	Where is the quote margin introducing direct speech positioned
				in relation to the quote content in natural texts? Is it always
				at the front? Can a quote margin ever come in the middle of a
				direct quote (as is sometimes found in \ili{English} literature)? Does
				it ever come at the end (like in journalistic \ili{English})? Do they
				have bipartite (bracketing) quote margins?

	\item	Are the complementisers used for speech act verbs also used
				with (some) verbs of cognition
				and perception? (\it{he spoke saying,
				he thought saying, he saw saying,
				{\ldots}})
				
	\item	Are repetition
				and tail-head linkage common features of narrative discourse
				in recorded texts? 
\end{itemize}